% Source-derived chapter 09: Quadratic forms
% Original source: standard-conjectures-abelian-fourfolds
\chapter{Quadratic forms}
\label{standard-conjectures-abelian-fourfolds-chapter-09}
\source{standard-conjectures-abelian-fourfolds}

\begin{remark}[Source section]
\label{standard-conjectures-abelian-fourfolds-chapter-09-source}
\source{standard-conjectures-abelian-fourfolds}
\source{standard-conjectures-abelian-fourfolds}
This chapter reproduces the corresponding section of the retrieved
LaTeX source, including its definitions, statements, examples, and
proofs. It has not yet been translated into Lean.
\notready
\end{remark}

% The source section title was: Quadratic forms
We recall here some classical facts on quadratic forms. They will allow us to reduce Theorem \ref{mainthm} to a $p$-adic question (Proposition \ref{padicreduction}). For simplicity, we will  work only in the context we will need later, namely with non-degenerate \nobreak{$\Q$-quadratic} forms of rank $2$.
In what follows $\Q_\nu$ denotes the completion of $\Q$ at the place $\nu$.
\begin{definition}
\source{standard-conjectures-abelian-fourfolds}
\notready\label{Hilbert}
\source{standard-conjectures-abelian-fourfolds}
Let $q$ be a $\Q$-quadratic form of rank $2$. Define $\varepsilon_\nu(q)$, the Hilbert symbol of $q$ at $\nu$, as  $+1$ if the equation
\[x^2-q(y,z)=0\]
has a nonzero solution in $x,y,z\in \Q_\nu$, and as $-1$ otherwise. Depending on the context we may write $\varepsilon_p(q)$ or $\varepsilon_\R(q)$.
\end{definition}

\begin{remark}
\source{standard-conjectures-abelian-fourfolds}
\notready\label{remarksignature}
\source{standard-conjectures-abelian-fourfolds}
Let $q$ be a $\Q$-quadratic form of rank $2$. It is positive definite if and only if its discriminant is positive and $\varepsilon_\R(q)=+1$
and it is negative definite if and only if its discriminant is positive and $\varepsilon_\R(q)=-1.$
\end{remark}

\begin{prop}[{\cite[\S 2.3]{cours}}]
\source{standard-conjectures-abelian-fourfolds}
\notready\label{localinvariant} 
\source{standard-conjectures-abelian-fourfolds}
Let $p$ be a prime number and $q_1$ and $q_2$   two non-degenerate $\Q$-quadratic forms of rank $2$. Then 
\[q_1 \otimes \Q_p \cong q_2 \otimes \Q_p\]
if and only if the discriminants of $q_1$ and $q_2$ coincide in $\Q_p^*/(\Q_p^*)^2$ and 
\[\varepsilon_p(q_1)=\varepsilon_p(q_2).\]
\end{prop}

\begin{theorem}[{\cite[\S 3.1]{cours}}]
\source{standard-conjectures-abelian-fourfolds}
\notready\label{productformula} 
\source{standard-conjectures-abelian-fourfolds}
Let $q$ be non-degenerate $\Q$-quadratic form of rank $2$ and $\varepsilon_\nu(q)$ be as in Definition \ref{Hilbert}. Then for all but finite places $\nu$ the equality $\varepsilon_\nu(q)=+1$ holds. Moreover, the following product formula running on all places holds
\[\prod_\nu  \varepsilon_\nu(q)= +1.\]
\end{theorem}

\begin{corollary}
\source{standard-conjectures-abelian-fourfolds}
\notready\label{signaturetop}
\source{standard-conjectures-abelian-fourfolds}
Let $q_1$ and $q_2$  be two non-degenerate $\Q$-quadratic forms of rank $2$ and let $p$ be a prime number. Suppose that, for all primes $\ell$ different from $p$, we have
\[q_1 \otimes \Q_\ell \cong q_2 \otimes \Q_\ell.\]
Then $q_1$ is positive definite if and only if one of the following two cases happens:
\begin{enumerate}
\item  The quadratic forms $q_1\otimes \Q_p$ and $q_2\otimes \Q_p$ are isomorphic and $q_2$ is positive definite.
\item  The quadratic forms $q_1\otimes \Q_p$ and $q_2\otimes \Q_p$ are not isomorphic and $q_2$ is negative definite.
\end{enumerate}
\end{corollary}
\begin{proof}

The $\ell$-adic hypothesis implies in particular that the discriminants of $q_1$ and $q_2$ coincide in $ \Q_\ell^*/(\Q_\ell^*)^2$ 
 for all   $\ell\neq p$. This implies that they coincide in $ \Q^*/(\Q^*)^2$
by \cite[Theorem 3 in 5.2]{Ireland}.

If the discriminants are negative, none of the conditions in the statement holds and the equivalence is clear. From now on we suppose that the discriminants are   positive. By Remark \ref{remarksignature}, $q_1$ is positive definite if and only $\varepsilon_\R(q_1)=+1.$




There are then two cases, namely $q_2$ is positive definite, respectively $q_2$ is negative definite. Again by Remark \ref{remarksignature} they are equivalent to $\varepsilon_\R(q_2)=+1$, respectively $\varepsilon_\R(q_2)=-1$. 

\


Now, Theorem \ref{productformula} implies that
$\prod_\nu  \varepsilon_\nu(q_1)=\prod_\nu  \varepsilon_\nu(q_2).$
Combining this with the $\ell$-adic isomorphisms we deduce
\[\varepsilon_\R(q_1)\varepsilon_p(q_1)=\varepsilon_\R(q_2)\varepsilon_p(q_2).\] This means that $q_1$ is positive definite if and only if
 \[\varepsilon_p(q_1)=\varepsilon_\R(q_2)\varepsilon_p(q_2).\]
 This relation is equivalent to the fact that one of the following two situations hold:
 \begin{enumerate}
\item   $q_2$ is positive definite and $\varepsilon_p(q_1)=\varepsilon_p(q_2).$
\item  $q_2$ is negative definite and $\varepsilon_p(q_1)\neq\varepsilon_p(q_2).$
\end{enumerate}
As the discriminant of $q_1$ and $q_2$ coincide, the equality $\varepsilon_p(q_1)=\varepsilon_p(q_2)$ is equivalent to the fact that $q_1\otimes \Q_p$ and $q_2\otimes \Q_p$ are isomorphic (Proposition \ref{localinvariant}).
\end{proof}


\begin{prop}
\source{standard-conjectures-abelian-fourfolds}
\notready\label{padicreduction}
\source{standard-conjectures-abelian-fourfolds}
Let us keep notation from Theorem \ref{mainthm}. Let $p$ be the characteristic of $k$ and $i$ be the unique non-negative integer such that the Hodge structure $V_B$ is of type $(i,-i)$ and  $(-i,i)$. Then, the quadratic form $q_Z$ is positive definite if and only if the  following   holds:
\begin{enumerate}
\item When $i$ is even: the quadratic forms $q_B\otimes \Q_p$ and $q_Z\otimes \Q_p$ are isomorphic.
\item  When $i$ is odd: the quadratic forms $q_B\otimes \Q_p$ and $q_Z\otimes \Q_p$ are not isomorphic.
\end{enumerate}
\end{prop}
\begin{proof}
By  Proposition \ref{tortura}(\ref{giorno}), we have that $q_Z \otimes \Q_\ell = R_\ell(q_{\vert{k}})$. By the comparison theorem, we have $q_B \otimes \Q_\ell = R_\ell(q_{\vert{\C}})$. Combining these equalities with smooth proper base change in $\ell$-adic cohomology   we deduce that \[q_B \otimes \Q_\ell = q_Z \otimes \Q_\ell.\]
On the other hand, by hypothesis, $q_B$ is a polarization for the Hodge structure $V_B$ hence it is positive definite if $i$ is even and negative definite if $i$ is odd. We can now conclude by applying Corollary \ref{signaturetop} to $q_1=q_Z$ and $q_2=q_B$.
\end{proof}
